\section{What's new}\label{whatsnew}

\subsection{Version 0.9.5}
\begin{compactitem}
\item Cross-platform support: the extension now builds and runs on
  Windows, macOS\,(11 Big Sur or newer, Apple Silicon only) and Linux.
\item \hyperref[units]{Extender support}: up to eight MCU-protocol units (one
  main plus up to seven extenders, i.e.\ up to 64 channels) are combined
  into one logical control surface. All channel-strip modes span the full
  channel count.
\item A lot of performance improvements make the extension more responsive,
  especially in larger projects.
\item The configuration dialog now configures every hardware unit (device
  type and MIDI ports); the former ``Icon QCon Pro X support'' checkbox has
  been replaced by a per-unit device type.
\item Improved plugin parameter mapping: discrete parameters (including
  their value names) are now detected automatically and mapped to the
  V-Pots.
\item The former ``2nd controller (Mackie Control B)'' workaround has been
  removed in favour of the built-in extender support.
\end{compactitem}

\subsection{Version 0.9}
\begin{compactitem}
\item \hyperref[prox]{Icon QCon Pro X support}
\item \hyperref[sendautomation]{Send/Receive automation}
\item Meterbridge improvements
\item New Plug-Mode Options (see Table 4)
\end{compactitem}

\subsection{Version 0.8}
\begin{compactitem}
\item \hyperref[buttonactions]{Button-emulating actions} (since v0.8.2)
\item \reaper folders can now be reflected on the MCU using the \foldermode
\item \anchors allows locking a \cs to a fixed track
\item \hyperref[fxpresets]{FX presets} (store and recall FX settings) including
  a Blind Test feature
\item \hyperref[localmaps]{FX Local maps} (different maps for each FX
  instance of the same FX)
\item improved syncing of the MCU state with the \reaper GUI
\item \hyperref[globalactions]{Global Actions} assignments can be named and
  shown in MCU display
\item \hyperref[quickjump]{Quick Jumps} added for faster navigation
\item improved \hyperref[tracknames]{Track Name} handling
\item FX favorites can be stored in \reaper project file
\end{compactitem}

